\def\mthreepure{1}
\documentclass[fontset=none]{ctexart}
\usepackage{qp-m3}
\begin{document}
\renewcommand{\qpzhangming}{第一章\quad 空间向量与立体几何}
\renewcommand{\qpjianming}{导学件·试产片4}
\keshi{第1课时\quad 空间向量概念（填空场）}
\begin{multicols}{2}
\emergencystretch=1em

\huaxing{课}{前}{预}{习}{知识导学\quad 素养初识}

\zsd{一}{空间向量的概念}

\tihao{1}{定义：在空间，具有\kongbai{}和\kongbai{}的量叫做空间向量；模为\(1\)的向量叫做\kongbai{}向量．}

\begin{ansblock}[A试-4-dx-01]
% ans:A试-4-dx-01 大小；方向；单位
\kbfill{大小|方向|单位}
\ansitem{1}{大小；方向；单位．}
\end{ansblock}

\tihao{2}{平面向量基本定理：同一平面内任一向量可由\kongbai{}的两个向量线性表示．}

\begin{ansblock}[A试-4-dx-02]
% ans:A试-4-dx-02 不共线
\kbfill{不共线}
\ansitem{2}{不共线．}
\end{ansblock}

\zhenhead{判断正误(正确的打\gou{},错误的打\cha{})}

\zhentib{(1)}{单位向量都相等．}[\cha]

\zhentib{(2)}{向量\(\overrightarrow{a}\)与\(-\overrightarrow{a}\)互为相反向量，它们的模相等．}[\gou]

\liubai[9mm]{此处书写}

\tihao{3}{共线向量定理：\(\overrightarrow{b}\)与非零\(\overrightarrow{a}\)共线\(\Leftrightarrow\)存在\kongbai{}个实数\(\lambda\)，使\(\overrightarrow{b}=\)\kongbai{}．}

\begin{ansblock}[A试-4-dx-03]
% ans:A试-4-dx-03 有且只有一；λa
\kbfill{有且只有一|\lambda\overrightarrow{a}}
\ansitem{3}{有且只有一；\(\lambda\overrightarrow{a}\)．}
\end{ansblock}

\end{multicols}

\end{document}
